\documentclass[12pt]{article}
\usepackage[utf8]{inputenc}
\usepackage[T1]{fontenc}
\usepackage{hyperref}
\title{Southern Marakwet Language Resources for K-12 Curriculum Development}
\author{Compiled by Claude Sonnet 5 (claude-sonnet-5) for Andr\'as Kornai}
\date{}
\begin{document}
\maketitle

\section{Introduction}
Southern Marakwet (ISO 639-3 tle) is a Kalenjin (Nilotic) language of Kenya. Note that the identification line in the task prompt gave the family as ``Bookkeeping'', which appears to be a data error. The sources found treat the language as Kalenjin. Global Recordings Network gives ``Southern Markweeta'' as an alternate name for Sambirir, and lists about 80,000 speakers. Joshua Project gives higher figures for Marakwet as a whole.

\section{Primary Educational Resources}
No primers, textbooks or school programmes were found.

\section{Linguistic and Reference Materials}
\begin{itemize}
\item Library of Congress authority record, which treats Endo and Sambirir as dialects and Marakwet as a cover term: \url{https://id.loc.gov/authorities/subjects/sh2004014786.html}
\item Marakwet people (background, with a reference list): \url{https://www.Wikipedia.org/wiki/Marakwet}
\end{itemize}

\section{Religious and Community Resources}
Global Recordings Network lists a New Testament (Faith Comes By Hearing) for the Sambirir and Cherang'any varieties: \url{https://www.globalrecordings.net/en/language/9569}. Joshua Project profile: \url{https://joshuaproject.net/index.php/people_groups/11760}

\section{Digital and Audio Resources}
The Global Recordings Network pages above list audio. A dictionary reportedly exists for Cherang'any, a related variety. Its orthography does not mark vowel quality or tone, and no link was found.

\section{Curriculum Development Context}
Materials aimed specifically at Southern Marakwet are scarce. The New Testament audio and the Cherang'any dictionary are the nearest resources, and they are for related varieties.
\end{document}